% GENERATED FILE. Edit canonical lessons, then run npm run generate:books.
% canonical-source-hash: fnv1a64:0217da8d33cebedd
% canonical-lessons: JA-C256-dou-yatte, JA-C256-mai, JA-C256-kippu-o-nimai, JA-C256-nihonjin, JA-W256-jin-kanji, JA-R256-which-way-to-the-station

\chapter{Which Way to the Station}
\label{ch:ja-which-way-to-the-station}
\hlchaptermodality{256}

\begin{chapteropening}
\textbf{By the end of this chapter:} \emph{I can ask the way (eki wa doko desu ka; kūkō made dō yatte ikimasu ka), follow directions (massugu itte, migi ni magatte kudasai), buy tickets with the counter mai, and ask when the next train leaves.}

Complete JA-R256-which-way-to-the-station: ask the way, buy two tickets and ask when the train leaves.
\end{chapteropening}

\section[\texorpdfstring{\ja{どこですか} … (doko desu ka …)}{どこですか … (doko desu ka …)}]{\ja{どこですか} / \ja{どうやって} (doko desu ka / dō yatte) — asking the way}

\label{lesson:JA-C256-dou-yatte}



\begin{quote}
Say \textquotedblleft{}where\textquotedblright{} and \textquotedblleft{}a station\textquotedblright{}. (\emph{doko}, \emph{eki}.) Then the \textbf{-\ja{て}} form of \textbf{\ja{いく}}. (\emph{itte}.)
\end{quote}

\subsection*{You'll want to know: \ja{どこですか}, \ja{どうやって}}
\begin{quote}
\textbf{\ja{えきは} \ja{どこですか。}}
\end{quote}

\begin{quote}
\emph{eki wa doko desu ka.} — \textquotedblleft{}Where is the station?\textquotedblright{}
\end{quote}

To ask how to get somewhere, \textbf{\ja{どうやって}} (\emph{dō yatte}) asks \textquotedblleft{}how\textquotedblright{}: \textbf{\ja{どう}}, \textquotedblleft{}how\textquotedblright{}, with \textbf{\ja{やって}}, the \textbf{-\ja{て}} form of \textbf{\ja{やる}}, \textquotedblleft{}to do\textquotedblright{}.

\begin{quote}
\textbf{\ja{くうこうまで} \ja{どうやって} \ja{いきますか。}}
\end{quote}

\begin{quote}
\emph{kūkō made dō yatte ikimasu ka.} — \textquotedblleft{}How do I get to the airport?\textquotedblright{}
\end{quote}

\begin{quote}
\textbf{\ja{ちかてつで} \ja{いきます。}}
\end{quote}

\begin{quote}
\emph{chikatetsu de ikimasu.} — \textquotedblleft{}You go by subway.\textquotedblright{}
\end{quote}

\subsection*{You'll want to know: \ja{いって、まがって} \ja{ください}}
Directions are \textbf{-\ja{て} \ja{ください}} requests, one after another, joined with the \textbf{-\ja{て}} form:

\begin{quote}
\textbf{\ja{まっすぐ} \ja{いって} \ja{ください。}}
\end{quote}

\begin{quote}
\emph{massugu itte kudasai.} — \textquotedblleft{}Please go straight on.\textquotedblright{}
\end{quote}

\begin{quote}
\textbf{\ja{まっすぐ} \ja{いって、みぎに} \ja{まがって} \ja{ください。}}
\end{quote}

\begin{quote}
\emph{massugu itte, migi ni magatte kudasai.} — \textquotedblleft{}Go straight on, then turn right.\textquotedblright{}
\end{quote}

\begin{quote}
\textbf{\ja{つぎの} \ja{かどを} \ja{ひだりに} \ja{まがって} \ja{ください。}}
\end{quote}

\begin{quote}
\emph{tsugi no kado o hidari ni magatte kudasai.} — \textquotedblleft{}Turn left at the next corner.\textquotedblright{}
\end{quote}

\begin{quote}
\textbf{\ja{ぎんこうの} \ja{まえに} \ja{あります。}}
\end{quote}

\begin{quote}
\emph{ginkō no mae ni arimasu.} — \textquotedblleft{}It is in front of the bank.\textquotedblright{}
\end{quote}

\subsection*{Your turn}
Say these aloud:

\begin{itemize}
\raggedright
  \item \textquotedblleft{}Where is the station?\textquotedblright{} — \emph{eki wa doko desu ka}
  \item \textquotedblleft{}Where is the bank?\textquotedblright{} — \emph{ginkō wa doko desu ka}
  \item \textquotedblleft{}Where is the hospital?\textquotedblright{} — \emph{byōin wa doko desu ka}
\end{itemize}

\begin{tcolorbox}[breakable,colback=teal!4,colframe=teal!35!black,title={Before you move on}]
\textquotedblleft{}How do I get there?\textquotedblright{} (\textbf{\ja{どうやって} \ja{いきますか}}.) \textquotedblleft{}Please turn right\textquotedblright{}? (\textbf{\ja{みぎに} \ja{まがって} \ja{ください}}.) \textquotedblleft{}Go straight on, then turn left\textquotedblright{}? (\textbf{\ja{まっすぐ} \ja{いって、ひだりに} \ja{まがって} \ja{ください}}.)
\end{tcolorbox}
\section[\texorpdfstring{\ja{まい} (mai)}{まい (mai)}]{\ja{まい} (mai) — the counter for flat things}

\label{lesson:JA-C256-mai}



\begin{quote}
Before the new word: say \textquotedblleft{}where is the station?\textquotedblright{} (\emph{eki wa doko desu ka}.) Then \textquotedblleft{}a ticket\textquotedblright{}. (\emph{kippu}.) People are counted with \textbf{\ja{にん}}; flat things have a counter of their own.
\end{quote}

\subsection*{You'll want to know: \ja{まい}}
\textbf{\ja{まい}} — \emph{mai} — the counter for flat things: tickets, stamps, sheets, plates. The number goes in front: \textbf{\ja{いちまい}} (\emph{ichimai}), one; \textbf{\ja{にまい}} (\emph{nimai}), two; \textbf{\ja{さんまい}} (\emph{sanmai}), three.

\begin{quote}
\textbf{\ja{きっぷを} \ja{にまい} \ja{ください。}}
\end{quote}

\begin{quote}
\emph{kippu o nimai kudasai.} — \textquotedblleft{}Two tickets, please.\textquotedblright{}
\end{quote}

\begin{quote}
\textbf{\ja{きってを} \ja{さんまい} \ja{ください。}}
\end{quote}

\begin{quote}
\emph{kitte o sanmai kudasai.} — \textquotedblleft{}Three stamps, please.\textquotedblright{}
\end{quote}

\begin{cousinweb}
\textbf{\ja{まい}} (\emph{mai}) is Sino-Japanese (kango): the on'yomi of a character KANJIDIC2 glosses \textquotedblleft{}sheet of …\textquotedblright{} and \textquotedblleft{}counter for flat thin objects or sheets\textquotedblright{}, a reading borrowed from Chinese (KANJIDIC2). JMdict lists it as a counter, \textquotedblleft{}for thin, flat objects (e.g. sheets of paper, plates, coins)\textquotedblright{} (JMdict).
\end{cousinweb}

\subsection*{Your turn}
Say these aloud:

\begin{itemize}
\raggedright
  \item \textquotedblleft{}One ticket, please\textquotedblright{} — \emph{kippu o ichimai kudasai}
  \item \textquotedblleft{}Two tickets, please\textquotedblright{} — \emph{kippu o nimai kudasai}
  \item \textquotedblleft{}Three tickets, please\textquotedblright{} — \emph{kippu o sanmai kudasai}
\end{itemize}

\begin{tcolorbox}[breakable,colback=teal!4,colframe=teal!35!black,title={Before you move on}]
Two tickets? (\textbf{\ja{にまい}}.) Three stamps? (\textbf{\ja{さんまい}}.) Two people? (\textbf{\ja{ふたり}}.)
\end{tcolorbox}
\section[\texorpdfstring{\ja{きっぷを} \ja{にまい} … (kippu o nimai …)}{きっぷを にまい … (kippu o nimai …)}]{\ja{きっぷを} \ja{にまい} \ja{ください} (kippu o nimai kudasai) — at the ticket window}

\label{lesson:JA-C256-kippu-o-nimai}



\begin{quote}
Say \textquotedblleft{}two tickets, please\textquotedblright{}. (\emph{kippu o nimai kudasai}.) Then \textquotedblleft{}a ticket window\textquotedblright{}, \textquotedblleft{}one way\textquotedblright{} and \textquotedblleft{}a round trip\textquotedblright{}. (\emph{madoguchi}, \emph{katamichi}, \emph{ōfuku}.)
\end{quote}

\subsection*{You'll want to know: \ja{〜}\ja{までの} \ja{きっぷ}}
\begin{quote}
\textbf{\ja{くうこうまでの} \ja{きっぷを} \ja{ください。}}
\end{quote}

\begin{quote}
\emph{kūkō made no kippu o kudasai.} — \textquotedblleft{}A ticket to the airport, please.\textquotedblright{}
\end{quote}

\textbf{\ja{まで}} says how far; \textbf{\ja{の}} joins it to the ticket: \textquotedblleft{}a ticket as far as the airport\textquotedblright{}. The number and \textbf{\ja{まい}} go before \textbf{\ja{ください}}:

\begin{quote}
\textbf{\ja{くうこうまでの} \ja{きっぷを} \ja{にまい} \ja{ください。}}
\end{quote}

\begin{quote}
\emph{kūkō made no kippu o nimai kudasai.} — \textquotedblleft{}Two tickets to the airport, please.\textquotedblright{}
\end{quote}

\begin{quote}
\textbf{\ja{かたみちですか。}}
\end{quote}

\begin{quote}
\emph{katamichi desu ka.} — \textquotedblleft{}One way?\textquotedblright{}
\end{quote}

\begin{quote}
\textbf{\ja{いいえ、おうふくを} \ja{にまい} \ja{ください。}}
\end{quote}

\begin{quote}
\emph{iie, ōfuku o nimai kudasai.} — \textquotedblleft{}No, two returns, please.\textquotedblright{}
\end{quote}

\subsection*{You'll want to know: \ja{なんじに} \ja{でますか}}
\textbf{\ja{なんじ}} (\emph{nanji}), \textquotedblleft{}what time\textquotedblright{}, is \textbf{\ja{なん}}, \textquotedblleft{}what\textquotedblright{}, with \textbf{\ja{じ}} (\emph{ji}), \textquotedblleft{}o'clock\textquotedblright{}. \textbf{\ja{でます}} (\emph{demasu}) is what a train does when it leaves:

\begin{quote}
\textbf{\ja{つぎの} \ja{でんしゃは} \ja{なんじに} \ja{でますか。}}
\end{quote}

\begin{quote}
\emph{tsugi no densha wa nanji ni demasu ka.} — \textquotedblleft{}What time does the next train leave?\textquotedblright{}
\end{quote}

\begin{quote}
\textbf{\ja{くじに} \ja{でます。}}
\end{quote}

\begin{quote}
\emph{kuji ni demasu.} — \textquotedblleft{}It leaves at nine.\textquotedblright{}
\end{quote}

\begin{quote}
\textbf{\ja{のりばは} \ja{どこですか。}}
\end{quote}

\begin{quote}
\emph{noriba wa doko desu ka.} — \textquotedblleft{}Where does it leave from?\textquotedblright{}
\end{quote}

\subsection*{Your turn}
Say these aloud:

\begin{itemize}
\raggedright
  \item \textquotedblleft{}It leaves at nine\textquotedblright{} — \emph{kuji ni demasu}
  \item \textquotedblleft{}It leaves at five\textquotedblright{} — \emph{goji ni demasu}
  \item \textquotedblleft{}It leaves at seven\textquotedblright{} — \emph{shichiji ni demasu}
\end{itemize}

\begin{tcolorbox}[breakable,colback=teal!4,colframe=teal!35!black,title={Before you move on}]
\textquotedblleft{}A ticket to the airport, please\textquotedblright{}? (\textbf{\ja{くうこうまでの} \ja{きっぷを} \ja{ください}}.) \textquotedblleft{}What time does it leave?\textquotedblright{} (\textbf{\ja{なんじに} \ja{でますか}}.) \textquotedblleft{}Where does it leave from?\textquotedblright{} (\textbf{\ja{のりばは} \ja{どこですか}}.)
\end{tcolorbox}
\section[\texorpdfstring{\ja{日本人} (nihonjin)}{日本人 (nihonjin)}]{\ja{日本人} (nihonjin) — a Japanese person}

\label{lesson:JA-C256-nihonjin}



\begin{quote}
Before the new word: say \textquotedblleft{}the Japanese language\textquotedblright{}. (\emph{nihongo}.) Then \textquotedblleft{}a friend\textquotedblright{}. (\emph{tomodachi}.)
\end{quote}

\subsection*{You'll want to know: \ja{日本人}}
\textbf{\ja{日本人}} — \emph{nihonjin} — \textquotedblleft{}a Japanese person\textquotedblright{}. It is \textbf{\ja{日本}} (\emph{nihon}), \textquotedblleft{}Japan\textquotedblright{}, the first part of \textbf{\ja{日本語}}, with \textbf{\ja{じん}} (\emph{jin}), \textquotedblleft{}person\textquotedblright{}, after it.

\begin{quote}
\textbf{\ja{ともだちは} \ja{日本人}\ja{です。}}
\end{quote}

\begin{quote}
\emph{tomodachi wa nihonjin desu.} — \textquotedblleft{}My friend is Japanese.\textquotedblright{}
\end{quote}

You write the first two signs. The third is new, read \textbf{\ja{じん}} (\emph{jin}) here; the next lesson writes it.

\begin{cousinweb}
\textbf{\ja{日本人}} (\emph{nihonjin}) is Sino-Japanese (kango): three Chinese characters read with their on'yomi, readings borrowed from Chinese. They are \emph{ni}, \textquotedblleft{}sun\textquotedblright{} (Middle Chinese \emph{nyit}), \emph{hon}, \textquotedblleft{}origin\textquotedblright{} (\emph{pwonX}) and \emph{jin}, \textquotedblleft{}person\textquotedblright{} (\emph{nyin}) (KANJIDIC2; Middle Chinese in Baxter's transcription). JMdict glosses it \textquotedblleft{}Japanese person\textquotedblright{} (JMdict).
\end{cousinweb}

\subsection*{Your turn}
Say these aloud:

\begin{itemize}
\raggedright
  \item \emph{nihonjin}
  \item \textquotedblleft{}My friend is Japanese\textquotedblright{} — \emph{tomodachi wa nihonjin desu}
  \item \textquotedblleft{}My teacher is Japanese\textquotedblright{} — \emph{sensei wa nihonjin desu}
\end{itemize}

\begin{tcolorbox}[breakable,colback=teal!4,colframe=teal!35!black,title={Before you move on}]
\textquotedblleft{}A Japanese person\textquotedblright{}? (\textbf{\ja{日本人}}, \emph{nihonjin}.) How is the new sign read here? (\textbf{\ja{じん}}, \emph{jin}.)
\end{tcolorbox}
\section[\texorpdfstring{\ja{人} (jin)}{人 (jin)}]{\ja{人} — a person, two strokes}

\label{lesson:JA-W256-jin-kanji}



\begin{quote}
Say \emph{nihonjin}, a Japanese person. You write \textbf{\ja{日}} and \textbf{\ja{本}}. Then call up \textbf{\ja{出}}: the long stem first.
\end{quote}

\subsection*{Script — the shape}
\begin{quote}
\textbf{\ja{人}} — \emph{jin}
\end{quote}

Its parts:

\begin{itemize}
\raggedright
  \item a stroke falling to the left
  \item a stroke falling to the right, leaning on it
\end{itemize}

Two strokes, one lift:

\begin{enumerate}
\raggedright
  \item draw the long falling stroke from the top, down to the left
  \item lift, start partway down the first stroke and draw the second down to the right
\end{enumerate}

The second stroke starts on the first, not at the top. Here it is read \textbf{\ja{じん}} (\emph{jin}), as in \textbf{\ja{日本人}}; this lesson asks only for that reading.

\subsection*{Your turn}
\begin{enumerate}
\raggedright
  \item Find \textbf{\ja{人}} in \textbf{\ja{日本人}} and say \emph{nihonjin}.
  \item Copy \textbf{\ja{人}} three times, saying \emph{jin} at each one.
  \item Find it once more, in \textbf{\ja{ともだちは} \ja{日本人}\ja{です}} (\emph{tomodachi wa nihonjin desu}), and say the line.
\end{enumerate}

\begin{tcolorbox}[breakable,colback=teal!4,colframe=teal!35!black,title={Before you move on}]
How many strokes? (\textbf{Two}.)

Source: \href{https://github.com/KanjiVG/kanjivg/blob/master/kanji/04eba.svg}{KanjiVG stroke-order data}: KanjiVG, kanji/04eba.svg, kanji \ja{人} (U+4EBA), 2 stroke paths kvg:04eba-s1 to s2 (Ulrich Apel and contributors, CC BY-SA 3.0). Only the order and direction are taken from it.
\end{tcolorbox}
\section[(review)]{Review — which way to the station}

\label{lesson:JA-R256-which-way-to-the-station}



\begin{quote}
Ask \textquotedblleft{}where is the station?\textquotedblright{} (\emph{eki wa doko desu ka}.) Then say \textquotedblleft{}two tickets, please\textquotedblright{}. (\emph{kippu o nimai kudasai}.)
\end{quote}

\subsection*{Your turn}
Say these aloud:

\begin{itemize}
\raggedright
  \item How do I get to the station?
  \item Turn left at the next corner.
  \item One ticket to the airport, please.
  \item What time does the next train leave?
  \item My friend is Japanese.
  \item Tomorrow is a holiday. So the bank is closed. — politely
\end{itemize}

\begin{tcolorbox}[breakable,colback=teal!4,colframe=teal!35!black,title={Before you move on}]
\textquotedblleft{}How\textquotedblright{}? (\textbf{\ja{どうやって}}.) Three stamps? (\textbf{\ja{さんまい}}.) \textquotedblleft{}What time\textquotedblright{}? (\textbf{\ja{なんじ}}.) \textquotedblleft{}A Japanese person\textquotedblright{}? (\textbf{\ja{日本人}}.)
\end{tcolorbox}
